
\documentclass[aspectratio=169]{beamer}

\mode<presentation>
\usetheme{Boadilla}
\definecolor{NTHU1}{RGB}{127,16,132}
\definecolor{NTHU2}{RGB}{228,210,231}
\definecolor{blue}{RGB}{30,90,205}
\definecolor{red}{RGB}{213,94,0}
\definecolor{green}{RGB}{0,128,0}
\definecolor{charcoalgray}{RGB}{108,104,104}
\setbeamercolor{title}{fg=NTHU1}
\setbeamercolor{frametitle}{fg=NTHU1}
\setbeamercolor{block title}{bg=NTHU1, fg=white}
\setbeamercolor{block body}{bg=NTHU2}
\setbeamercolor{structure}{fg=NTHU1}
\setbeamercolor{item projected}{fg=white}
\setbeamercolor{item}{fg=NTHU1}
\setbeamercolor{subitem}{fg=NTHU1}
\setbeamercolor{section in toc}{fg=NTHU1}
\setbeamercolor{description item}{fg=NTHU1}
\setbeamercolor{caption name}{fg=NTHU1}
\setbeamercolor{button}{bg=NTHU1, fg=white}
\setbeamercolor{caption name}{fg=NTHU1}
\usepackage{graphics}
\usepackage{tikz}
\usepackage{amsmath}
\usepackage{bbm}
\usetikzlibrary{decorations.pathreplacing}
\usepackage{geometry}
\usepackage{booktabs}
\usepackage{bookmark}
\usepackage{multirow, makecell}
\usepackage{float}
\usepackage{fancyvrb}
\usepackage{caption}
\captionsetup[figure]{singlelinecheck = false, format= hang, justification=centering}
\captionsetup[subfigure]{singlelinecheck = false, format= hang, justification=raggedright}
\usepackage{subcaption}
\usepackage{adjustbox}
\usepackage{hyperref}
\usepackage{threeparttable}
\usepackage{appendixnumberbeamer} 
\usepackage[T1]{fontenc}
\usepackage[default]{lato} 
\usepackage{smartdiagram}
\smartdiagramset{
  back arrow disabled = true,
  module minimum width=1cm,
  module x sep = 3,
  uniform arrow color=true,
}
\usepackage{xcolor}
\usepackage{colortbl}
	\definecolor{light-gray}{gray}{0.99}

\newenvironment{wideitemize}{\itemize\addtolength{\itemsep}{10pt}}{\enditemize}
\newenvironment{wideenumerate}{\enumerate\addtolength{\itemsep}{10pt}}{\endenumerate}
\newenvironment{widedescription}{\description\addtolength{\itemsep}{10pt}}{\enddescription}
\hypersetup{
colorlinks=true,
linkcolor=NTHU1,
filecolor=green, 
urlcolor=blue,
}
\beamertemplatenavigationsymbolsempty
\setbeamercolor{author in head/foot}{bg=white, fg=NTHU1}
\setbeamercolor{title in head/foot}{bg=white, fg=NTHU1}
\setbeamercolor{date in head/foot}{bg=white, fg=NTHU1}
\setbeamercolor{section in head/foot}{bg=white, fg=NTHU1}
\setbeamercolor{headline}{bg=white}
\setbeamertemplate{footline}{
    \leavevmode%
    \hbox{%
        \begin{beamercolorbox}[wd=.333333\paperwidth,ht=2.25ex,dp=1ex,center]{date in head/foot}%
            \usebeamerfont{date in head/foot}%\insertdate
        \end{beamercolorbox}%
        \begin{beamercolorbox}[wd=.333333\paperwidth,ht=2.25ex,dp=1ex,center]{title in head/foot}%
            \usebeamerfont{title in head/foot}\insertshorttitle
        \end{beamercolorbox}%
        \begin{beamercolorbox}[wd=.333333\paperwidth,ht=2.25ex,dp=1ex,right]{date in head/foot}%
            \usebeamerfont{date in head/foot}\insertframenumber{}\hspace*{2em}
        \end{beamercolorbox}}%
        \vskip0pt%
    }
    %% May need to script the introduction!!!!!!
%\setbeamercolor{page number in head/foot}{fg=NTHU1}
\setbeamertemplate{section in toc}[sections numbered]
\setbeamertemplate{subsection in toc}{\leavevmode\leftskip=3em\rlap{\hskip-1.75em\inserttocsectionnumber.\inserttocsubsectionnumber}
\inserttocsubsection\par}
\setbeamerfont{subsection in toc}{size=\footnotesize}
%\setbeamertemplate{headline}{%
  %\begin{beamercolorbox}[ht=5.5ex]{section in head/foot}
    %\vskip2pt\insertnavigation{0.33\paperwidth}\vskip2pt
  %\end{beamercolorbox}%
%}
\newenvironment{transitionframe}{\setbeamercolor{background canvas}{bg=white}\setbeamertemplate{footline}{} \begin{frame}[noframenumbering]}{\end{frame}}


\makeatletter
\let\@@magyar@captionfix\relax
\makeatother


\title[]{Public Economics and Development} % Change this regularly
\subtitle[]{Week 10: Types of public services: Cash vs. In-kind, Conditional vs. Unconditional}
\author[Lee]{Seung-hun Lee }
\institute[]{Taipei School of Economics\\ National Tsing Hua University}

\date[]{April 28th, 2026}

\begin{document}
\begin{transitionframe}
\titlepage
\end{transitionframe}


%%% Color slides for section headers: Use for colloquium version (The ones bewteen \iffals and \fi)

\section{Overview setup of the topic}

\begin{frame}
\frametitle{Widespread usage of Cash and In-kind transfers}
\begin{figure}
  \centering
  \includegraphics[keepaspectratio, width=0.9\textwidth]{fig1.png}
\end{figure}
\begin{footnotesize}
Source: Gentilli, U. (2023) ``Why does In-kind Assistance Persist when Evidence Favors Cash Transfers?'', Brookings Institution
\end{footnotesize}
\end{frame}



\begin{frame}
\frametitle{Cash vs.\ in-kind transfers}
\begin{columns}
\begin{column}{0.48\textwidth}
  \textbf{Case for cash}
  \begin{itemize}\setlength\itemsep{4pt}
    \item Recipients know their own needs best -- cash respects consumer sovereignty
    \item Less administrative cost (no procurement, storage, or distribution)
    \item Recipients prefer the cash equivalent and reallocate accordingly than in-kind
  \end{itemize}
\end{column}
\begin{column}{0.48\textwidth}
  \textbf{Case for in-kind}
  \begin{itemize}\setlength\itemsep{4pt}
    \item Paternalism: recipients may not spend cash on socially desired goods
    \item Self-targeting: in-kind goods are less attractive to the non-poor
  \item Promote consumption of goods generating positive externality

  \end{itemize}
\end{column}
\end{columns}
\medskip
\begin{block}{Key question}
\begin{itemize}
\item[$\to$] Respect recipient preferences, or promote socially desired consumption?
\item[$\to$] Targeting or insurance value of in-kind transfers $>$ Higher cost relative to cash?
\end{itemize}
\end{block}
\end{frame}

\begin{frame}
\frametitle{Conditional cash transfers: Progresa in Mexico}
Mexico's Progresa (1997-2019) as the pioneering large-scale CCT
\begin{itemize}\setlength\itemsep{3pt}
\item Targets poor rural households via proxy-means testing
\item Transfers conditional on school attendance (85\%+) and health clinic visits
\item By 2004, covered 5 million households --- around a quarter of Mexico's population
\item Enrollment, child health , household consumption $\uparrow$ (effect wiped out post program)
\item Spawned a global wave: Used in 2 countries in 1997 $\to$ 63 LMICs in 2016

\end{itemize}\medskip
Why Progresa became a blueprint
\begin{itemize}\setlength\itemsep{3pt}
\item One of the first programs to be rigorously evaluated via randomized rollout
\item Demonstrated simultaneous gains in enrollment, health, and consumption
\item Spawned a large literature debating on optimal design
 \begin{itemize}\setlength\itemsep{3pt}
\item How much do conditions matter? 
\item Who benefits most? 
\item What is the optimal CCT/UCT mix?
\end{itemize}
\end{itemize}
\end{frame}


\begin{frame}
\frametitle{Unconditional cash transfers: GiveDirectly in Kenya}
An NGO-driven large-scale unconditional cash transfer
\begin{itemize}\setlength\itemsep{3pt}
\item Rolled out in poor households in rural Kenya via M-Pesa mobile money
\item Eligibility based on a simple means test (thatched roof as poverty proxy) 
\item Average transfer of \$709 PPP, delivered as lump sum or monthly installments
\begin{itemize}\setlength\itemsep{3pt}
\item Roughly two years of per capita expenditure, 
\end{itemize}
\item Became the benchmark against which other transfer modalities are compared
\end{itemize}\medskip
What the evidence shows 
\begin{itemize}\setlength\itemsep{3pt}
\item Monthly household consumption from \$158 to \$193 PPP nine months after transfer
\item Significant gains in assets, food security, and psychological well-being
\item No evidence of wasteful spending (alcohol, tobacco)
\item Then, is there a need to condition benefits at all?
\end{itemize}
\end{frame}




\begin{frame}
\frametitle{The CCT vs.\ UCT debate}
\begin{columns}
\begin{column}{0.48\textwidth}
  \textbf{Case for CCT}
  \begin{itemize}\setlength\itemsep{4pt}
    \item Correct underinvestment in human capital due to misperceived returns 
\item Imposing a condition deters the non-needy from applying
    \item Taxpayers accept conditioned redistribution more readily
  \end{itemize}
\end{column}
\begin{column}{0.48\textwidth}
  \textbf{Case for UCT}
  \begin{itemize}\setlength\itemsep{4pt}
\item More autonomy for households
\item Conditions distort choices: households may comply not because it is optimal but because it is required
\item Simpler and cheaper to administer (no need to monitor compliance)
  \end{itemize}
\end{column}
\end{columns}
\medskip
\begin{block}{Key breaking point}
\begin{itemize}
\item[$\to$] If the goal is welfare improvement, do conditions add beyond the income transfer?
\item[$\to$] If the goal is behavior change, are conditions the most efficient way to achieve it?
\end{itemize}
\end{block}
\end{frame}

\begin{transitionframe}
\frametitle{Roadmap for today}
\tableofcontents
\end{transitionframe}








\section{Cash vs In-kind Transfers}
\subsection{Hidrobo, Hoddinott, Peterman, Margolies, Moreira (2014): Cash, Food or Vouchers? Evidence from Randomized Experiment in Northern Ecuador}
\begin{transitionframe}
\begin{center}
  \begin{huge}
    \textcolor{NTHU1}{%
Hidrobo, Hoddinott, Peterman, Margolies, Moreira (2014):\\ Cash, Food or Vouchers? Evidence from Randomized Experiment in Northern Ecuador
    }
  \end{huge}
\end{center}
\end{transitionframe}


\begin{frame}
\frametitle{Q: How does form of assistance matter for welfare outcomes? } 
Debate over the form of assistance has long history
\begin{itemize}\setlength\itemsep{3pt}
\item Proponents of cash
\begin{itemize}\setlength\itemsep{3pt}
\item Autonomy: Allow beneficiaries to use transfers as they see fit
\item Less stigma: Less visible compared to in-kind transfers
\end{itemize}
\item Proponents of in-kind transfers
\begin{itemize}\setlength\itemsep{3pt}
\item Targeting: Only those who truly need it will take-up these benefits
\item Paternalism: In-kind transfers facilitate changes in a particular behavior
\end{itemize}
\end{itemize}\medskip
This paper uses RCT to comparatively evaulate the effects of cash and in-kind transfers
\begin{itemize}\setlength\itemsep{3pt}
\item Effect of cash, food vouchers, and food transfers on food consumption
\item Track quantity and quality (calory intake, dietary diversity)
\item Also provides analysis on cost-effectiveness across different programs
 \end{itemize}
\end{frame}

\begin{frame}
\frametitle{Program design} 
Background: WFP assistance to poor households and Colombia refugees in Ecuador
\begin{itemize}\setlength\itemsep{3pt}
\item 6 Monthly transfers of cash, food, and vouchers $\simeq\$40$ (May - Oct 2011)
\item Eligibility determined by proxy means test based on surveys
\item Cash transfers sent to ATM cards and did not expire
\item Vouchers redeemable for nutritionally-approved foods at markets (nontransferrable)
\item Food transfers gave rice, oil, lentiloes, canned sardines
\end{itemize}\medskip
Some characteristics of the recipients
\begin{itemize}\setlength\itemsep{3pt}
\item Mostly women: Original intent was to increase women's role in household decisions
\item Mostly poor: Dirt floors, rarely own TV/computer/automobiles
\item Has 1-2 children (mostly aged between 6-15)
\item All received same lecture on significance of nutrition
\end{itemize}
\end{frame}




\begin{frame}
\frametitle{Research design} 
Experiment details
\begin{itemize}\setlength\itemsep{3pt}
\item Randomization at neighborhood  level
\item Neighborhoods randomized into one of the four treatments
\item Thus, treatment arms are comparable at baseline 
\item Regression based on OLS (Analysis of Covariance)
\[
Y_{hj,1}=\beta_f\text{food}_j+\beta_c\text{cash}_j\beta_v\text{voucher}_j + \gamma Y_{hj,0} + \text{StrataFE} +\text{NeighborFE}+e_{hj}
\]
\end{itemize}\medskip
Additional data obtained
\begin{itemize}\setlength\itemsep{3pt}
\item Following quality measures are used
\begin{itemize}\setlength\itemsep{3pt}
\item Dietary diversity (DDI): Sum of no. of distinct food items consumed last week (1-40)
\item Household dietary diversity (HDDS): no. of distinct food groups consumed (1-12) 
\item Food cosumption score (FCS): $\simeq$ HDDS, different weights per food group (35 as cutoff)
\end{itemize}
\item Caloric intake constructed from Nutritional Database for Standard Reference (USDA)
\end{itemize}
\end{frame}


\begin{frame}
\frametitle{Overall improvements across treatment groups}
\begin{figure}
  \centering
  \includegraphics[keepaspectratio, width=0.7\textwidth]{hhpmm2014_f1.png}
\end{figure}
\end{frame}

\begin{frame}
\frametitle{All treatments show increases in quantity of consumption}
\begin{figure}
  \centering
  \includegraphics[keepaspectratio, width=1\textwidth]{hhpmm2014_t1.png}
\end{figure}
\end{frame}

\begin{frame}
\frametitle{Food transfer $\to$ Calorie intake $\uparrow$ and Voucher $\to$ DDI}
\begin{figure}
  \centering
  \includegraphics[keepaspectratio, width=1\textwidth]{hhpmm2014_t2.png}
\end{figure}
\end{frame}


\begin{frame}
\frametitle{Cost-effectiveness and remainig discussion} 
Estimates of cost-effectiveness
\begin{itemize}\setlength\itemsep{3pt}
\item Costs: WFP's accounting ledgers
\begin{itemize}\setlength\itemsep{3pt}
\item Modality-specific staff costs, debit card production, storage facilities etc
\end{itemize}
\item Costs: Food transfers $>$ Cash $\simeq$ Vouchers
\begin{itemize}\setlength\itemsep{3pt}
\item For each transfer: \$11.46, \$ 2.99, and \$3.27 each
\item For raising certain outcomes,  Food transfers $>$ Cash $>$ Vouchers
\end{itemize}
\item Recipient preference for cash: 9\% wish not to receive further help vs. 30\% in others
\end{itemize}\medskip
Discussion
\begin{itemize}\setlength\itemsep{3pt}
\item Choosing the best mode of assistance depends on specific objectives
\item Simply improve welfare? Cash is preferable
\item Promote particular behaviors? In-kind might be better
\item Considering recipient preference? In-kind may not be best suited
\end{itemize}
\end{frame}




\subsection{Cunha, De Giorgi, Jayachandran (2019): The Price Effects of Cash vs In-Kind Transfers}
\begin{transitionframe}
\begin{center}
  \begin{huge}
    \textcolor{NTHU1}{%
Cunha, De Giorgi, Jayachandran (2019):\\ The Price Effects of Cash vs In-Kind Transfers
    }
  \end{huge}
\end{center}
\end{transitionframe}


\begin{frame}
\frametitle{Q: How does different anti-poverty program affect prices?}
How do in-kind and cash transfers affect market circumstances?
\begin{itemize}\setlength\itemsep{3pt}
\item Cash transfers: Increase demand for normal goods 
\begin{itemize}\setlength\itemsep{3pt}
\item Think of how expansionary policy affects price levels
\end{itemize}
\item In-kind: Increase demand for normal goods, but also increase supply
\item Potentially have different effects depending on competition structure
\begin{itemize}\setlength\itemsep{3pt}
\item Perfect competition: Shift wealth between buyers and sellers
\item Imperfect competition: Correct inefficiencies by reducing high prices
\end{itemize}
\end{itemize}\medskip
This paper tests the presence of these price effects in rural Mexico
\begin{itemize}\setlength\itemsep{3pt}
\item Compare between status quo, cash transfers, and in-kind transfers
\item Estimates changes in prices following inflow of support
\item Disaggregates size of price effects across villages with different development
\end{itemize}
\end{frame}

\begin{frame}
\frametitle{Conceptual framework: Perfect comptition} 
\begin{columns}
\begin{column}{0.45\textwidth}
  \begin{figure}
    \centering
  \includegraphics[keepaspectratio, width=\textwidth]{cdj2019_f1.png}
  \end{figure}
\end{column}
\begin{column}{0.55\textwidth}
  Cash transfers
  \begin{itemize}\setlength\itemsep{3pt}
    \item Demand curve shift right ($\because$ income $\uparrow$)
    \item Hence, price $\uparrow$
  \end{itemize}\medskip
  In-kind transfer (equivalent value)
  \begin{itemize}\setlength\itemsep{3pt}
    \item Demand curve shift right ($\because$ income $\uparrow$)
    \item Some of the demand provided by govt
    \item Residual demand for local suppliers shift left
    \item Hence, price $\downarrow$ 
    \begin{itemize}\setlength\itemsep{3pt}
      \item In-kind good demand do not rise drastically
      \item Unless it is a luxury good
    \end{itemize}
  \end{itemize}\medskip
\end{column}
\end{columns}
\end{frame}


\begin{frame}
\frametitle{Conceptual framework: Imperfect comptition} 
\begin{columns}
\begin{column}{0.45\textwidth}
  \begin{figure}
    \centering
  \includegraphics[keepaspectratio, width=\textwidth]{cdj2019_f2.png}
  \end{figure}
\end{column}
\begin{column}{0.55\textwidth}
  Assume $N$ firm Cournot-Nash competition
  \begin{itemize}\setlength\itemsep{3pt}
    \item Constant MC at $c$, Demand $p=d-Q$
    \item Equilibrium $p=\frac{d+Nc}{N+1}$
  \end{itemize}\medskip
  Cash transfers
  \begin{itemize}\setlength\itemsep{3pt}
    \item Demand curve shift right ($\because$ income $\uparrow$)
    \item Hence, price $\uparrow$
  \end{itemize}\medskip
  In-kind transfer (equivalent value)
  \begin{itemize}\setlength\itemsep{3pt}
    \item Demand curve shift right ($\because$ income $\uparrow$)
    \item Residual demand for local suppliers shift left
    \item Hence, price $\downarrow$ 
    \item Room to correct excessively high price
  \end{itemize}\medskip
\end{column}
\end{columns}
\end{frame}


\begin{frame}
\frametitle{Experiment design and context} 
Programa de Apoyo Alimentario (PAL)
\begin{itemize}\setlength\itemsep{3pt}
\item 5,000 small ($<2,500$ ppl) highly marginalized villages
\item Monthly allotment of basic food items (rice, beans etc) and supplements (tuna etc)
\end{itemize}\medskip
Experiment design
\begin{itemize}\setlength\itemsep{3pt}
\item Randomly selected 208 villages
\item Assigned to In-kind, Cash transfer, and control group
\item Value of cash is 150 MXP: 18\% of food expense \& 8\% of income per month
\item Value of in-kind is roughly 206 MXP, but w/ assumptions, have similar value to cash
\begin{itemize}\setlength\itemsep{3pt}
\item 116 MXP is consumed
\item 35 MXP of extra consumption of transferred goods, 55 MXP not consumed 
\item The latter two discounted to 2/3 $\to$ in-kind value of 146 MXP
\end{itemize}
\item Uses household and store surveys to obtain data on price
\end{itemize}
\end{frame}


\begin{frame}
\frametitle{Price in in-kind fell relative to cash villages} 
\begin{figure}
  \centering
  \includegraphics[keepaspectratio, width=1\textwidth]{cdj2019_t1.png}
\end{figure}
\end{frame}



\begin{frame}
\frametitle{Price differences larger in less developed villages} 
\begin{figure}
  \centering
  \includegraphics[keepaspectratio, width=0.9\textwidth]{cdj2019_t2.png}
\end{figure}
\end{frame}


\begin{frame}
\frametitle{Other findings and key takeaways}
Other findings
\begin{itemize}\setlength\itemsep{3pt}
\item The price effects are actually persistent
\begin{itemize}\setlength\itemsep{3pt}
\item Villages treated early had price effects even when later-treated villages began treatement
\end{itemize}
  \item Purchasing power in developed villages do not change much
  \item Larger price effect in poor villages due to less market integration and competition
  \item Local supplier profit increases in cash transfer villages (no change in in-kind village)

\end{itemize}\medskip
Discussions
\begin{itemize}\setlength\itemsep{3pt}
\item In-kind transfers can correct prices set excessively high due to imperfect competition
\item But if government is inefficient, this effect could be limited
\item If so, cash transfers with policies to promte supply-side competition may be preferred
\item Choice of cash vs. in-kind depends on policy goals and market circumstances
\end{itemize}
\end{frame}





\subsection{Gardenne, Norris, Singhal, Sukhtankar (2024): In-kind Transfers as Insurance}
\begin{transitionframe}
\begin{center}
  \begin{huge}
    \textcolor{NTHU1}{%
Gardenne, Norris, Singhal, Sukhtankar (2024): \\In-kind Transfers as Insurance
    }
  \end{huge}
\end{center}
\end{transitionframe}



\begin{frame}
\frametitle{Q: How does in-kind transfers insurce against price volatility?} 
Low-income households are vulnerable to price variations
\begin{itemize}\setlength\itemsep{3pt}
\item Chages in price of key consumption goods affect stability of their budget
\item Particularly for developing countries: Market integration $\downarrow$ and weather shocks
\item If households only receive cash transfers, they cannot insure against such shock
\item The insurace aspect of transfers, not jus transfers, is understudied
\end{itemize}\medskip
This paper sheds light on the insurace aspect of in-kind transfers
\begin{itemize}\setlength\itemsep{3pt}
  \item Conceptual framework of price-indexed transfers equalizing marginal utility of income
  \item In-kind transfers can provide additional insurance value on top of redistrubution
  \begin{itemize}\setlength\itemsep{3pt}
  \item Especially when price-indexing is difficult due to lack of high-frequency data
\end{itemize}
\item Empirically tests the idea using calorie intake as proxy for marginal utility of income
\begin{itemize}\setlength\itemsep{3pt}
  \item Using a flagship in-kind transfer program in India (Public Distribution System, PDS)
\end{itemize}
\end{itemize}
\end{frame}



\begin{frame}
\frametitle{Conceptual Framework} 
Optimal insurance policy
\begin{itemize}\setlength\itemsep{3pt}
\item Consider households consuming several goods, where one of its price $p_j$ varies 
\item It maximizes indirect utility $v(p, y)$ ($y$: income)
\item Insurance: Price-dependent transfers $x(p_j)$
\item Optimal insurance equates marginal utility of income ($\mu$) in all states of the world
\[
v_y(p, y+x(p_j)) = \mu
\]\vspace{-20pt}
\item If $v_y$ (MU of income) increase in price, then $x(p_j)$ increases in price
\end{itemize}\medskip
Role of in-kind transfers (vs. benchmark cash transfer)
\begin{itemize}\setlength\itemsep{3pt}
\item Same goods in fixed quantity vs same expected monetary value across all states
\item $p_j\uparrow\to$ Market value of an in-kind transfer $\uparrow\to$  purchasing power $\uparrow$ 
\item In-kind transfers have insurance value (price-indexed transfers) -- not fixed transfers
\item This hinges on whether additional income more valuable when prices are high

\end{itemize}\medskip
\end{frame}

\begin{frame}
  \frametitle{Empirical testing}
India and PDS system
\begin{itemize}\setlength\itemsep{3pt}
\item Prices are volatile and many fail to meet caloric requirements
\item PDS provides rice and other staple foods at below-market prices
\end{itemize}\medskip
Transferring theory to data
\begin{itemize}\setlength\itemsep{3pt}
\item MU of income proxied with indicator for failing to meet calorie requirements
\begin{itemize}\setlength\itemsep{3pt}
\item Assume: Likelihood of failing to meet requirments $\uparrow$ associated with MU of income $\uparrow$
\item If unlikely to meet requirement, additional income valuable for meeting this requirement
\end{itemize}
\item Relevant price and calorie intake data from National Sample Survey (NSS, 2003-12) 
\item Also control for household characteristics (obtained from NSS)
\item Uses the following regression to estimate relationship between price and caloric intake
\[
c_{it}=\beta \text{Riceprice}_{-i,t}+X_{it}\lambda+\delta_{da}+\theta_t + \phi_{\text{round}}+e_{it}
\]\vspace{-20pt}
where leave-out price is used to net out influence of own-taste
\end{itemize}
\end{frame}

\begin{frame}
  \frametitle{More likely to not meet requirements with higher prices}
  \begin{figure}
    \centering
    \includegraphics[keepaspectratio,width=0.9\textwidth]{gnss2024_t1.png}
  \end{figure}
\end{frame}

\begin{frame}
  \frametitle{Particularly for poor rural households}
  \begin{figure}
    \centering
    \includegraphics[keepaspectratio,width=0.9\textwidth]{gnss2024_t2.png}
  \end{figure}
\end{frame}


\begin{frame}
  \frametitle{Also reduces sensitivity of meeting requirements to market prices}
  \begin{figure}
    \centering
    \includegraphics[keepaspectratio,width=0.67\textwidth]{gnss2024_t3.png}
  \end{figure}
\end{frame}


\begin{frame}
\frametitle{Key takeaways and discussion} 
Further findings
\begin{itemize}\setlength\itemsep{3pt}
\item Value of in-kind that would ensure meeting caloric requirements
\begin{itemize}\setlength\itemsep{3pt}
  \item Find $x$ s.t. $-0.260 + 0.177 \times x =0$
  \item This $x$ is 147 Rupees, 1.5 times the average nonzero transfer
\end{itemize}
\item Thus, PDS implemented in the sample period did contribute to substantial insurnace
\end{itemize}\medskip
Policy implications
\begin{itemize}\setlength\itemsep{3pt}
  \item So does in-kind transfers dominate cash transfers? Not necessarily
  \item Need to consider implementation costs (As discussed in the previous paper)
  \item Nevertheless, insurance value provided by in-kind could be worthwhile for recipients
  \item Especially when recipients are exposed to price risk that affects their welfare
  \end{itemize}
\end{frame}

\section{Conditional vs. Unconditional transfers}
\subsection{Baird, McIntosh, Özler (2011): Cash or Condition? Evidence from a Cash Transfer Experiment}
\begin{transitionframe}
\begin{center}
  \begin{huge}
    \textcolor{NTHU1}{%
Baird, McIntosh, Özler (2011):\\ Cash or Condition? Evidence from a Cash Transfer Experiment
    }
  \end{huge}
\end{center}
\end{transitionframe}



\begin{frame}
\frametitle{Q: How to conditioning benefits affect recipient behaviors?} 
Relative merits between conditioning (CCT) and not conditioning (UCT) cash transfers?
\begin{itemize}\setlength\itemsep{3pt}
 \item Evidence that both CCT and UCT do affect behaviors across different settings
 \item Proponents of CCT 
\begin{itemize}\setlength\itemsep{3pt}
\item Correct for underinvestment in education and health by imposing conditions
\item Make transfers politically acceptable to the non-poor
\end{itemize}
\item Proponents of UCT
\begin{itemize}\setlength\itemsep{3pt}
\item Marginal contribution of conditioning unknown
\item Conditioning may be difficult to implement in dedveloping countries
\end{itemize}
\end{itemize}\medskip
This paper examines marginal contribution of conditioning cash transfers
\begin{itemize}\setlength\itemsep{3pt}
  \item Cash transfers to households with school-age girls in Malawi
  \item RCT where some receive CCT, others are under UCT, and remaining  receive none
  \item Examine behaviors that CCT conditions (education) 
  \item But also investigate others (human capital, underage marriage / pregnancy)
  \end{itemize}
\end{frame}



\begin{frame}
\frametitle{Experimental setting and design}
Zomba District in Malawi
\begin{itemize}\setlength\itemsep{3pt}
    \item One of the poorest in the world (2008 GNI per capita is \$760)
    \item Enrollment to secondary school is low (24\%)
    \item 176 Enumeration Areas spread across urban, near rural, far rural areas
    \item Surveyed to identify never-married females aged 13-22
\end{itemize}\medskip

Treatment assignment
\begin{itemize}\setlength\itemsep{3pt}
    \item CCT arm (T1, 46 EAs): Monthly transfer conditional on regular school attendance
    \begin{itemize}\setlength\itemsep{3pt}
    \item  Transfer amount determined by random: \$4 - \$10 per month
    \item Separate payment to schoolgirls (\$1-\$5) to cover school fees
\end{itemize}
\item UCT arm (T2, 27 EAs): No school requirement, same financial offer
\item Remaining EAs (88) are in control group
\item Track attendence, performance, marriage, pregnancy for `08 and `09
\item Potential concern: Selective attrition (especially for control group) 
  \end{itemize}
\end{frame}




\begin{frame}
\frametitle{No selective attrition!}
\begin{figure}
  \centering
  \includegraphics[keepaspectratio, width=1\textwidth]{bmo2011_t1.png}
\end{figure}
\end{frame}

\begin{frame}
\frametitle{Higher enrollment in CCT (based on teacher reporits)}
\begin{figure}
  \centering
  \includegraphics[keepaspectratio, width=1\textwidth]{bmo2011_t2.png}
\end{figure}
\end{frame}

\begin{frame}
\frametitle{Better educational performance from CCT}
\begin{figure}
  \centering
  \includegraphics[keepaspectratio, width=0.9\textwidth]{bmo2011_t3.png}
\end{figure}
\end{frame}


\begin{frame}
\frametitle{Less early marriage and pregnancy in UCT}
\begin{figure}
  \centering
  \includegraphics[keepaspectratio, width=0.9\textwidth]{bmo2011_t4.png}
\end{figure}
\end{frame}


\begin{frame}
\frametitle{Key findings and discussion} 
Other key findings
\begin{itemize}\setlength\itemsep{3pt}
  \item UCT impacts on marriage and pregnancy entirely driven by girls who dropped out
  \begin{itemize}\setlength\itemsep{3pt}
  \item Effects of UCT on these outcomes null among those enrolled
  \end{itemize}
  \item Authors do a follow-up (paper in 2019 at \textit{Journal of Development Economics})
\begin{itemize}\setlength\itemsep{3pt}
  \item Besides sustained enrollment, effects of CCT vanish
  \item Health gains of UCT besides height-for-age vanish
  \end{itemize}
  \item Effects of CCT similar across different amounts of transfers
\end{itemize}\medskip

Discussion
\begin{itemize}\setlength\itemsep{3pt}
    \item CCT can be cost-effective in raising educational attainment 
    \item But UCTs still can improve other outcomes (health, delayed marriage)
    \item CCT could improve outcomes that hinge on complying to conditions (e.g. test scores)
    \item UCT effective if many noncompliers change behaviors following income support
\end{itemize}

\end{frame}




\subsection{Bergstrom, Dodds (2021): The Targeting Benefit of Conditional Cash Transfer}
\begin{transitionframe}
\begin{center}
  \begin{huge}
    \textcolor{NTHU1}{%
Bergstrom, Dodds (2021):\\ The Targeting Benefit of Conditional Cash Transfer
    }
  \end{huge}
\end{center}
\end{transitionframe}

\begin{frame}
\frametitle{Q: Does conditional transfers have a targeting \textit{benefit}?}
Major opposition to conditional cash transfer (CCT): Targeting cost
\begin{itemize}\setlength\itemsep{3pt}
\item CCTs are conditional on investments in children's human capital (schooling or health)
\item Low-income households find these conditions costly $\to$ excluded (targeting cost)
\item Thus, many prefer UCTs (no targeting costs) to alleviate poverty
\end{itemize}\medskip
This paper argues that CCTs have targeting benefits that could offset targeting costs
\begin{itemize}\setlength\itemsep{3pt}
\item Sending a chilren: lumpy investment where child income is the opportunity cost
\item CCTs direct money to households forgoing child income (targeting benefit)
\item Provides a theoretical framework capturing targeting benefits and costs
\item Proposes a optimal trade-off between CCT and UCT with estimates from rural Mexico
\end{itemize}
\end{frame}


\begin{frame}
\frametitle{Theoretical framework}
Set-up of the household optimziation
\begin{itemize}\setlength\itemsep{3pt}
\item Parents decide whether to send children to work (for income $y_c$) or to school
\begin{itemize}
\item Schooling HH forgo $y_c$ but receive higher income when children become adults
\end{itemize}
\item All receive UCT $t_u$, schooling HH receive $t_c$ as well (along with parental income $y$)
\item Cannot borrow against child's future income (Otherwise, always consumption more)
\item Thus, they solve ($s$: schooling, $\mu$: child's ability)
\[
\max_{s\in\{0,1\}} u(c) + v(\mu, s) \ \ \text{(s.t.) } c = y + y_c(1-s)+t_u + t_c s
\]
\end{itemize}\medskip
Social planner problem
\begin{itemize}\setlength\itemsep{3pt}
\item Utilitarian: Maximizes sum of social welfares of schooling \& non-schooling HH
\item Distribute fixed budget to predetermined set of households in UCT or CCT
\item Can observe $s$ (not $y$, or $\mu$) - so that CCTs goes to schooling households ($s=1$)
 \end{itemize}
\end{frame}

\begin{frame}
\frametitle{Key takeaways from the framework}
Conclusions: $(1+\frac{\partial t_u}{\partial t_c})\overline{u_c(c(1))}S+\frac{\partial t_u}{\partial t_c}\overline{u_c(c(0))}(1-S)=0$
\begin{itemize}\setlength\itemsep{3pt}
\item Targeting benefit: Net welfare to schooling HH when $t_c$ increases
\item Targeting cost: Net welfare loss from decreasing $t_u$ due to rise in $t_c$
\item Key component: MU of consumption higher in schooling households!
\begin{itemize}\setlength\itemsep{3pt}
  \item Especially if $y_c$ is large and parental income is largely equal
  \item May not hold if schooling households are rich to begin with (most developed countries)
\end{itemize}
\item In such case, better to increase budget to CCT 
\end{itemize}\medskip
How to capture from empirically observable objects (sufficient statistics)?
\begin{itemize}\setlength\itemsep{3pt}
\item $S$: Share of schooling households (obtainable from data)
\item $\overline{u_c(c(s))}$: Consumption distribution across households + curvature of utility functions
\item $\frac{\partial t_u}{\partial t_c}$: Budget trade-off (Not observable on its own)
\begin{itemize}\setlength\itemsep{3pt}
  \item Need income effect of UCT (I), price effect of CCT (P), and $S$
  \item I tracks how $S$ changes with UCT, P tracks how $S$ changes with CCT
\end{itemize}
\end{itemize}
\end{frame}


\begin{frame}
  \frametitle{Taking the theory to the data}
  Progresa program in Mexico
\begin{itemize}\setlength\itemsep{3pt}
\item Largest component: nutritional subsidy and education grant
\item 320 treated localities (1998 --), 186 control localities (2000 --)
\item Transfers conditional on children attending school 85\%+ of days (grades 3--9)
\end{itemize}\medskip
Identifying income effects ($I$)
\begin{itemize}\setlength\itemsep{3pt}
\item Use variation in transfers to siblings under age 12
\item Enrollment below age 12 is near 100\% (practically a UCT)
\item Variation in sibling composition provides exogenous variation in size of $t_u$
\end{itemize}\medskip
Identifying price effects ($P$)
\begin{itemize}\setlength\itemsep{3pt}
\item Introduction of CCTs was randomized at the locality level
\item Transfer schedule varies by grade and gender, providing variation in size of $t_c$
\end{itemize}
\end{frame}

\begin{frame}
  \frametitle{CCTs can narrow the consumption gap}
  \begin{figure}
    \centering
    \includegraphics[keepaspectratio, width=0.7\textwidth]{bd2021_f1.png}
  \end{figure}
\end{frame}

\begin{frame}
  \frametitle{$P=0.008$, $I=0.004$}
  \begin{figure}
    \centering
    \includegraphics[keepaspectratio, width=1\textwidth]{bd2021_t1.png}
  \end{figure}
\begin{itemize}
  \item 1 Peso increase in CCT $\to$ raise schooling by 0.8 p.p
  \item 1 Peso increase in UCT $\to$ raise schooling by 0.4 p.p
\end{itemize}
\end{frame}


\begin{frame}
  \frametitle{Optimal CCT/UCT mix: quantitative results}
  How much of the budget should go to a CCT on targeting grounds alone?
\begin{itemize}\setlength\itemsep{3pt}
\item Large forgone $y^c$ and limited parental income inequality $\to$ TB is substantial
\item From Progresa: Targeting benefit is 79\% of the targeting cost 
\item TB/TC $<$ 1 $\Rightarrow$ current CCT share is slightly too high; optimal to shift some to UCT
\item CCT/UCT mix that equalizes this: 33\% should go to CCT
\end{itemize}\medskip
Takeaway for policy design
\begin{itemize}\setlength\itemsep{3pt}
\item The targeting benefit is a quantitatively important advantage of CCTs \item It should not be ignored when designing cash transfer programs
\item CCT vs. UCT is not a binary choice: an optimal program allocates some budget to each
\item CCT preferred when $c(1) < c(0)$ and enrollment is below socially optimal
\item UCT preferred when $c(1) > c(0)$ and enrollment is near optimal
\end{itemize}
\end{frame}

\begin{frame}
\frametitle{Unanswered questions}
Program design aspects
\begin{itemize}\setlength\itemsep{3pt}
\item Optimal program size: Needs vs fiscal sustainability and general equilbrium effects?
\item Amount of support -- Universal amount or vary with circumstances?
\end{itemize}\medskip
Behavioral and long-run effects
\begin{itemize}\setlength\itemsep{3pt}
\item Do transfers build long-run resilience or create dependency? 
\item What about life-cycle effects on labor supply, savings, and investment?
\item Many papers try to look at graduation programs to study this

\end{itemize}\medskip
Room for further research
\begin{itemize}\setlength\itemsep{3pt}
\item Role of digital money: How can it close the gaps?
\begin{itemize}\setlength\itemsep{3pt}
  \item Individuals can send money back to families as remittance
  \item Potential insurance role: Value of transfers adjust to real-time valuation
\end{itemize}

\end{itemize}
\end{frame}



%%%%%%%%%%%
\end{document}
